% Generated draft: requires mathematical and source review.
\documentclass[11pt]{article}
\usepackage[margin=25mm]{geometry}
\usepackage{fontspec}
\setmainfont{Latin Modern Roman}
\usepackage{amsmath,amssymb,mathtools,unicode-math}
\setmathfont{Latin Modern Math}
\usepackage{enumitem,needspace,xurl,hyperref}
\setlength{\parindent}{0pt}
\setlength{\parskip}{5pt}
\setlength{\emergencystretch}{3em}
\newcommand{\N}{\mathbb{N}}
\newcommand{\Z}{\mathbb{Z}}
\newcommand{\Q}{\mathbb{Q}}
\newcommand{\R}{\mathbb{R}}
\newcommand{\C}{\mathbb{C}}
\newcommand{\F}{\mathbb{F}}
\newcommand{\A}{\mathbb{A}}
\newcommand{\PP}{\mathbb{P}}
\newcommand{\E}{\mathbb{E}}
\newcommand{\eps}{\varepsilon}
\newcommand{\dd}{\,\mathrm{d}}
\begin{document}
\section*{UnsolvedMath ID: 20001157}
\textbf{Source problem code:} catalogue entry
\subsection{Definitions and mathematical statement}
A finite simple graph is balanced here when it has no odd cycle and no induced cycle of length congruent to \(2\pmod4\); holes have length at least four. The Conforti--Rao source conjecture is that every balanced graph with at least one edge has an edge \(e\) for which \(G-e\) remains balanced. The edge-existence qualification is explicit because the literal assertion cannot hold for an edgeless graph. The source attributes the contribution to Gerard Cornuejols.

Call an edge safe if deleting it preserves balancedness. The aid reports the following exact certificate for an unsafe edge \(e=uv\): there are two internally vertex-disjoint \(u\)-\(v\) paths \(P_1,P_2\), each of length congruent to \(3\pmod4\), such that the induced subgraph on their union is exactly
\[
 P_1\cup P_2\cup\{uv\}.
\]
Thus \(uv\) is the one-edge branch of an induced theta, or equivalently the unique chord of the forbidden-length cycle formed by the other two branches. The criterion presumes \(G\) itself is balanced.

In particular every edge incident with a vertex of current degree at most two is reported safe. Repeatedly deleting such edges preserves balancedness and leaves exactly the edge set of the original \(3\)-core, with possibly isolated vertices outside it. The \(3\)-core is the unique maximal induced subgraph of minimum degree at least three. Therefore every balanced \(2\)-degenerate graph, meaning every nonempty induced subgraph has a vertex of degree at most two, has an ordering of all its edges in which every successive deletion preserves balancedness. This reduction does not assert that all remaining core edges are unsafe and does not solve the general conjecture, which the aid reports still open.

The record also contains the separate subsection ``K.1 Balanced circulants.'' For integers \(k,m\ge1\), set \(N=4km\) and define \(G(k,m)\) on \(\Z_N\), with an edge between distinct \(i,j\) exactly when
\[
 i-j+4mt\equiv\pm1\pmod N\quad\text{for some }t\in\Z,
\]
equivalently \(i-j\equiv\pm1\pmod{4m}\). For \(k=5,m=2\), the source gives differences
\(\{1,7,9,15,17,23,25,31,33,39\}\) modulo \(40\). It states that \(G(k,m)\) is balanced, that its holes have lengths only \(4\) and \(4m\), and that every cycle of length \(2\pmod4\) has at least two chords. Hence it reports \(G(k,m)-e\) balanced for every edge \(e\), a stronger family-specific deletion property. These circulant claims are retained as source statements, not inferred as a general solution. The trailing phrase ``in agreement with'' is incomplete and supplies no additional assertion.
\subsection{Short English statement}
The source conjectures that a nontrivial balanced graph always has a safely deletable edge. The supplied aid characterizes unsafe edges by an induced modular theta and proves safe deletion outside the three-core, while the source's separate circulant family permits every single-edge deletion.
\subsection{Sources}
\begin{itemize}
\item Local UnsolvedMath registry, permanent ID \(20001157\), source code \texttt{catalogue entry}; \texttt{aim-combinatorics-notes.json}, \texttt{notes[281]}, tag \texttt{conjecture}.
\item AIM workshop ``The Perfect Graph Conjecture,'' Conjecture 2, Conforti--Rao, contributed by Gerard Cornuejols, followed by ``K.1 Balanced circulants'': \url{https://aimath.org/WWN/perfectgraph/perfectgraph.pdf}. Statement status \texttt{exact}. No PDF was accessed for this draft.
\item No original remarks, literature field, full external conjecture citation, or authoritative proof artifact for the candidate unsafe-edge and core statements is supplied.
\item Research attempt 1: \texttt{partial\_result}, problem status at run \texttt{open}, rigor label \texttt{proved}, novelty confidence low; registry status \texttt{partially\_solved}, classification \texttt{PARTIAL-PROGRESS}, suggested difficulty L4. Created and updated 14 August 2026. This machine-generated aid is not peer reviewed; independently verify before citation.
\item Ulam AI / UnsolvedMath Contributors; CC BY 4.0.
\end{itemize}
\end{document}
